\documentclass[11pt,a4paper]{article}
\usepackage[margin=23mm]{geometry}
\usepackage{amsmath,amssymb,booktabs,array,longtable,hyperref}
\hypersetup{hidelinks}
\emergencystretch=3em
\newcommand{\Lc}{\textbf{(c)}\ }
\newcommand{\Ld}{\textbf{(d)}\ }
\newcommand{\dd}{\mathrm d}
\newcommand{\ccheck}{\texttt{onebody\_\allowbreak scheme\_\allowbreak continuum.py}}
\newcommand{\dcheck}{\texttt{onebody\_\allowbreak scheme\_\allowbreak dictionary.py}}
\newcommand{\echeck}{\texttt{onebody\_\allowbreak scheme\_\allowbreak edge.py}}
\title{The one-body factor of four:\\an audit and a lattice-matched continuum prescription}
\author{Research note for Gustav Delius}
\date{4 November 2026}
\begin{document}
\maketitle

\begin{abstract}
\Lc The failed assertion in \texttt{audit\_onebody.py} is an unjustified
accuracy claim about finite-$h$ extrapolations, not a failure of the limiting
edge identity. Its maximum residual is $2.2546734882\,10^{-4}$, not below
$3\,10^{-6}$. Independent inverse-Laplace calculations and fresh finite-sea
solves diagnose the error.

\Lc A more consequential omission is the continuum vacuum elastic
renormalisation. In the specified good-equation prescription its sound-speed
coefficient is $-9/(4\pi)$; under the \emph{same formal continuation} as the
soliton trace it is $+9/(4\pi)$ for the bad equation. Toda's harmonic vacuum
gives $+3/(4\pi)$. Matching these independently calculated vacuum coefficients
requires the finite Hamiltonian density
$\Delta\mathcal H=-3\hbar_B u^2/(2\pi)$, and changes the continued soliton
coefficient from $2/\pi$ to exactly $1/(2\pi)$. Both bad prescriptions then
have the same sound-subtracted coefficient $-1/(4\pi)$.

\Lc Conclusion \textbf{(A), for the leading coefficient in the stipulated
continued trace}: a precise lattice-matched prescription exists, fixed by
vacuum matching rather than a soliton fit. This is an effective-field-theory
finite matching term, not just a different Wick contraction of $u^3$.
\Ld It does not construct the physical constrained bad-equation determinant,
prove equivalence of complete quantum theories, or make the real stable-band
trace agree with Toda. Those stronger identifications remain open.
\end{abstract}

\section{Scope, conventions and checks}
Labels are (a) proved in the literature, (b) conjectured or heuristic in the
literature, (c) derived or checked here, and (d) unestablished identification
or guess. Numerical (c) means evidence at the stated precision, not a rigorous
error bound. No external article, abstract, DOI or arXiv entry is cited.
The relevant internal sources are identified in Section~\ref{sec:sources}.
All scripts are in \texttt{phase2/scripts/}; each writes only its own
\texttt{onebody\_scheme\_*.json} and ends with
\texttt{ALL CHECKS PASSED}. We leave all existing files and git state alone.

\Lc (\ccheck, \dcheck) Use the stipulated potentials and Hamiltonians,
with $u=\phi_x$, $\pi=\phi_t$, and $\Omega=\int\dd\pi\wedge\dd\phi$:
\begin{align}
 \mathcal H_g&=\tfrac12\pi^2+\tfrac12u^2-u^3+\tfrac12u_x^2,
 &v_g&=\sqrt{1-k^2},\\
 \mathcal H_b&=\tfrac12\pi^2+\tfrac12u^2+u^3-\tfrac12u_x^2,
 &v_b&=\sqrt{1+k^2}.
\end{align}
Here $k=\sin\alpha$ on the good branch and $k=\sinh\beta$ on the bad
branch. We display right movers; parity supplies $|P|$ for left movers.
For $u_0=(k^2/2)\operatorname{sech}^2(kx/2)$ the checked integrals are
\begin{equation}\label{eq:moments}
 \int u_0\dd x=2k,\qquad \int u_0^2\dd x=\frac23k^3,
 \qquad \int u_0^3\dd x=\frac4{15}k^5,\qquad
 P_0=v\int u_0^2\dd x.
\end{equation}
Throughout $F$ and $W$ denote \emph{coefficients} of $\hbar_B$ in the
fixed-$P$ soliton energy and homogeneous vacuum energy density, respectively.
The notation $s_1$ means $c_{\rm sound}=1+\hbar_Bs_1+\cdots$ in fixed
classical Boussinesq units. We test candidate schemes; none is adopted as the
PI's preferred quantum theory.

\section{Assessment of audit\_onebody}\label{sec:assessment}
\subsection{The failed assertion: limiting identity versus extrapolation}
\Lc (\echeck) The original script was executed unchanged and failed at
line 383, as reported. Its comparison is
\begin{equation}\label{eq:identity}
 C_v(y)=s_{1,T}+\frac y2\{u'(y)-1\},\qquad
 s_{1,T}=\frac3{4\pi}.
\end{equation}
This is an $h\to0$ \emph{edge-limit} identity. The script substitutes
independently extrapolated entries from \texttt{onebody\_toda\_results.json}
and \texttt{frame\_analysis.json}, with finite-$h$ truncation errors.
The input solves are double precision: evaluating the comparison at 65
digits does not improve their extrapolation accuracy. The central fits use
only three $h$ values for a quadratic intercept. At large $y$, the finite
amplitude $\kappa^2=hy/b$ is not uniformly small on those grids.

\begin{center}\small
\begin{tabular}{rrrr}\toprule
$y$ & original signed residual & combined fit spread & finer quartic residual\\\midrule
8 & $-7.5273\,10^{-7}$ & $2.3259\,10^{-5}$ & ---\\
16 & $-5.0951\,10^{-6}$ & $8.4061\,10^{-5}$ & $-1.6466\,10^{-9}$\\
32 & $-3.4752\,10^{-5}$ & $2.9805\,10^{-4}$ & $-4.1939\,10^{-8}$\\
64 & $-2.2547\,10^{-4}$ & $1.1013\,10^{-3}$ & $-9.3801\,10^{-7}$\\\bottomrule
\end{tabular}
\end{center}
\Lc The residual is $C_v-s_1-y(u'-1)/2$. Every original residual lies
inside its recorded combined sensitivity spread. Fitting the finer saved
finite-sea rows independently, with a quartic through the last seven
$h$ values down to $.0025$, recovers the original $3\,10^{-6}$ tolerance.
At $y=64$, a quintic gives $-2.2681\,10^{-7}$; a cubic gives
$-5.5779\,10^{-6}$ and still fails that tolerance. These are model
sensitivities, not certified bounds. We do not relax or edit the old check.

\paragraph{Independent calculation, not only refitting.}
\Lc (\echeck) Factoring the half-line kernel at nonzero Fourier frequency
gives the checked product and inverse transforms
\begin{align}
 g(s)&=s^{-1/2}\Gamma\!\left(1+\frac{s}{2\pi}\right)
       \exp\!\left[-\frac{s}{2\pi}\log\frac{s}{2\pi e}\right],\\
 2\pi g(iw)g(-iw)&=\frac{2\pi}{1-e^{-|w|}},\quad w\ne0,\\
 \mathcal L[F](s)&=\sqrt\pi\,g(s)/s^2,\quad
 \mathcal L[f](s)=\sqrt\pi\,g(s)/s,\label{eq:transforms}\\
 \mathcal L[J](s)&=\sqrt\pi\,g(s)/s^3,\\
 F(u(y))&=\frac43y^{3/2},\quad f=F',\quad J'=F.
\end{align}
The normalization is fixed by $F\sim4u^{3/2}/3$, not by the target
$C_T$. Fresh Talbot inversions and de Hoog contour checks at 45 digits
agree within $10^{-35}$ in all three transforms. Differentiation and
integration of the dressed momentum give
\begin{equation}\label{eq:edgeobservables}
 C_v=\frac3{4\pi}+\frac{3F}{4f}-\frac y2,
 \qquad C_E=\frac3{4\pi}+\frac{3J}{4F}-\frac{3y}{10},
 \qquad u'=\frac{2\sqrt y}{f}.
\end{equation}
Thus the fresh velocity results are
\begin{center}\small
\begin{tabular}{rrrr}\toprule
$y$ & fresh $C_v$ & old intercept minus fresh & finer intercept minus fresh\\\midrule
16 & $0.159293610890713$ & $-4.5672\,10^{-6}$ & $-1.5142\,10^{-9}$\\
64 & $0.159189963972784$ & $-2.0347\,10^{-4}$ & $-8.7287\,10^{-7}$\\\bottomrule
\end{tabular}
\end{center}
\Lc Separate fresh composite-Gauss finite-sea solutions at
$h=.02,.01,.005$ use $800,1600,3200$ nodes, with a 4000-node
refinement at $.005$. They reproduce the finer saved
velocities within $10^{-7}$; the refinement changes the tested quantities
by less than $10^{-11}$. Their \emph{finite-$h$} residual decreases
towards zero. These computations use neither the old solver nor $C_T$ as
an input. This distinguishes truncation error from an incorrect limiting
identity.

\subsection{Six dictionary factors}
\Lc (\dcheck) All six leading classical conversions have factor one.
The following checks make clear what this assertion does and does not fix.
Write $a$ for mean Toda spacing and $b=e^{-a/2}$, with reference coordinates
\begin{equation}\label{eq:dictionary}
 x_B=2\sqrt3\,n,\quad t_B=2\sqrt3 b\,t_T,
 \quad \phi_B=-U_T/\sqrt3,\quad
 S_T=3bS_B,\quad \hbar_T=3b\hbar_B.
\end{equation}
Here $U_T$ is displacement; $u$ remains Boussinesq strain. The Taylor
expansion of the Toda action through cubic strain and the first dispersive
term, and the symplectic form, independently give the same factor $3b$.

\begin{enumerate}
\item \textbf{Planck constant.} The action factor fixes $\hbar_T=3b\hbar_B$;
there is no extra factor four in the semiclassical exponential.
\item \textbf{Momentum definitions.} The canonical soliton label momentum
is $P_n=4b(\kappa\cosh\kappa-\sinh\kappa)$, not the total particle
momentum $\Pi_T=2b\sinh\kappa$. The physical-frame canonical momentum is
$P_n/a$. Their different small-$\kappa$ orders prohibit interchanging them.
\item \textbf{Dressed momentum.} Independently integrate
$p+(a\pi)^{-1}\int_0^\pi\Theta_0(p-2b\cos t)\dd t$, where
$\Theta_0(d)=2d(\log d-1)$ and $p=2b\cosh\kappa$.
It equals $P_n/a$ to better than $10^{-53}$ at four test points.
At finite $h$, the checked origin is $P(q)=\pi h/a$; use
$\overline P=P(p)-P(q)$.
\item \textbf{Energy and its origin.} Use $\overline E=E(p)-E(q)$ and
$E_{T,0}=2b^2(\sinh\kappa\cosh\kappa-\kappa)$. The conversion is
\begin{equation}\label{eq:units}
 P_B=\frac{a\overline P}{6\sqrt3 b},\quad
 E_B=\frac{\overline E}{6\sqrt3 b^2},\quad
 \frac{\Delta E_B}{\hbar_B P_B}
 =\frac3a\frac{\Delta\overline E}{h\overline P}.
\end{equation}
With $\kappa=\sqrt3\beta$, the leading momentum, energy, and fifth-order
excess above the classical sound line all match. An energy-origin shift
cannot change a velocity coefficient.
\item \textbf{Mass and width.} $M_B=2k$, $M_T=-2\kappa$, and
$M_T/(-\sqrt3M_B)\to1$. Eliminating the width in favour of the canonical
momentum removes a mere width reparametrisation.
\item \textbf{Scattering phase.} Independent evaluation of the two exact
Hirota polynomials in overtaking and head-on channels gives
$\delta_T/(3b\delta_b)\to1$. At amplitude scale $10^{-5}$ the errors are
$1.11\,10^{-10}$ and $6.49\,10^{-11}$. Division by the correspondingly
scaled $\hbar$ leaves factor one.
\end{enumerate}
\Lc These checks exclude the proposed \emph{multiplicative classical}
explanations. They do not determine an additive order-$\hbar$ elastic
counterterm, a quantum frame convention, or a projected determinant domain.
The audit's conclusion that the dictionary rules out a scheme explanation
does not follow.

\subsection{Toda decomposition and qualifications to its proof}
\Lc (\dcheck, \echeck) Fresh 60-digit finite-sea solves check the
endpoint identity $D(q)V(q)=\pi h/a$, the edge momentum origin, and
\begin{equation}\label{eq:thermo}
 v_s=-a\partial_aq\big|_h
 =\frac a2(q-h\partial_hq\big|_a),\qquad
 V(p)=\partial_a(aP(p))\big|_{p,h}.
\end{equation}
At $h=.8,q=1$, independent $q$ and $h$ differences with step $10^{-15}$
and 48 Gauss nodes satisfy the specified $10^{-27}$ tolerance. Endpoint
products and derivatives are not evaluated by substituting these identities
into themselves.

\Lc (\ccheck) The Toda sound coefficient has an independent vacuum
derivation. The harmonic frequency at spacing $a$ is
$\omega_T(K)=2b|\sin(K/2)|$, so the ground-state energy per particle is
\begin{equation}\label{eq:todavac}
 e(a,h)=b^2+h\frac1{4\pi}\int_0^{2\pi}\omega_T(K)\dd K+O(h^2)
       =b^2+\frac{2hb}{\pi}+O(h^2).
\end{equation}
With unit particle mass, the long-wave sound speed is
$a\sqrt{\partial_a^2 e}=ab+ah/(4\pi)+O(h^2)$; the dictionary gives
$s_{1,T}=3/(4\pi)$ in Boussinesq units.

\Lc (\echeck) Expanding \eqref{eq:transforms} at small $s$ and inverting
the first logarithm gives
\begin{equation}\label{eq:edgelog}
 u(y)=y-\frac{\log y}{2\pi}
       +\frac{1-\log(8\pi)}{2\pi}+O(y^{-1}\log^m y).
\end{equation}
Insertion into \eqref{eq:edgeobservables} yields, for the stipulated nested
edge limit,
\begin{equation}\label{eq:todac}
 C_T=\frac3{4\pi}-\frac1{4\pi}=\frac1{2\pi}.
\end{equation}
Thus the decomposition is confirmed; ``universal'' here refers to this
matched acoustic edge problem, not arbitrary UV completions or quantum
frames. \Ld The audit's claimed proof of commutation with a fixed
nonzero-$P$ semiclassical limit substitutes the large-$y$ expansion into a
different limit without a uniform remainder estimate. Its formal result is
compatible with \eqref{eq:todac}, but that substitution is not a proof of
limit interchange.

\paragraph{Additional errors in the audit text.}
\Lc (\dcheck, \echeck) Theorem 3 prints $b=\pi/a$; the correct relation
is $b=e^{-a/2}$. Its intermediate expression for $(V/D-v_s)/h$ incorrectly
subtracts the classical $ay/6$ term: before subtracting classical soliton
velocity the edge expression is $ay u'(y)/6$. The displayed final
\eqref{eq:identity} is correct. If the inverse edge series contains
$e_0/y^2$, integrating with $\dd P\propto\sqrt y\dd y$ gives the energy
tail $3e_0/y^2$, not the printed $3e_0/(2y^2)$. None of these text errors
changes the leading constant derived above.

\subsection{What the short-wave claim establishes}
\Lc (\ccheck) The conditional good trace and its formal continuation are
\begin{align}
 F_g(k)&=\frac{\sqrt3}{\pi}(\tfrac34-k^2)
            \operatorname{artanh}(2k/\sqrt3)-\frac{3k}{2\pi},\\
 F_{\rm ac}(k)&=-iF_g(ik)
 =\frac{\sqrt3}{\pi}(\tfrac34+k^2)
            \arctan(2k/\sqrt3)-\frac{3k}{2\pi}.
 \label{eq:F}
\end{align}
They give $C_g=-2/\pi$ and $C_{\rm ac}=2/\pi$. For the three algebraic
roots $\varphi=\alpha,\alpha\pm2\pi/3$, the coefficient of the logarithm
in $e-vp$ is proportional to
\begin{equation}\label{eq:weights}
 w(\varphi)=2\sin\varphi(\cos\alpha-\cos\varphi),\quad
 w(\alpha)=0,\quad
 w(\alpha\pm2\pi/3)=\pm2\sqrt3(\tfrac34-\sin^2\alpha).
\end{equation}
This exactly confirms the \emph{logarithmic} light-root cancellation.
It does not prove that the spectral integral has support only near
$|q|=\sqrt3/2$. In particular the elementary term of the light root in
the displayed charge sum is $-k/\pi$, not zero; all elementary root terms
together give $-3k/(2\pi)$. Allocation among roots, endpoint and tadpole
is representation dependent. The accurate statement is that the final
nonpolynomial term in \eqref{eq:F} comes from the two heavy algebraic roots.
The literal ``100\% short-wave modes'' assertion overinterprets this algebra.

\Lc (\ccheck) The audit also incorrectly describes the real-band
resonance. At $k\to0$, for a right mover and $q_*=-\sqrt3/2$,
\begin{equation}\label{eq:resonance}
 \omega(q_*)=\frac{\sqrt3}{4},\quad
 \omega'(q_*)=1,\quad
 \omega(q_*)-q_*=\frac{3\sqrt3}{4}\ne0.
\end{equation}
It is a stationary co-moving frequency ($\omega'-v\to0$), not a zero
co-moving frequency $\omega=-vq$. The positive-$q$ group derivative there
is $-1$ and reverses sign at negative $q$; the audit's $-1/2$ is wrong.
The maximum of $\omega(q)$ on the positive band is at $1/\sqrt2$, not
$\sqrt3/2$. The latter lies on the descending part, still inside the
stable band.

\section{Vacuum calculation and a precise matching prescription}
\subsection{Good continuum: the missing acoustic renormalisation}
\Lc (\ccheck) Put a constant strain $u=U+\eta_x$ on the good vacuum.
The quadratic frequency at fixed reference coordinates is
\begin{equation}\label{eq:homogeneous}
 \omega_g(q;U)=|q|\sqrt{q^2+1-6U}.
\end{equation}
Its zero-point energy coefficient at a fixed wave-number cutoff is
\begin{equation}\label{eq:bareW}
 W_{g,\Lambda}^{\rm bare}(U)
 =\frac1{2\pi}\int_0^\Lambda q\sqrt{q^2+1-6U}\dd q
 =\frac{(\Lambda^2+1-6U)^{3/2}-(1-6U)^{3/2}}{6\pi}.
\end{equation}
Subtract the vacuum constant and exactly its first derivative at $U=0$.
The latter is the specified normal-ordering tadpole, since
$(W_{g,\Lambda}^{\rm bare})'(0)=-3(\sqrt{1+\Lambda^2}-1)/(2\pi)$.
The remaining limit is
\begin{equation}\label{eq:Wg}
 W_g(U)=-\frac{(1-6U)^{3/2}-1+9U}{6\pi},\qquad
 W_g''(0)=-\frac9{2\pi}.
\end{equation}
The convergent curvature is independently checked by
\begin{equation}
 W_g''(0)=-\frac9{2\pi}\int_0^\infty\frac{q\dd q}{(1+q^2)^{3/2}}
         =-\frac9{2\pi}.
\end{equation}
Two direct subtracted quadratures at finite $\Lambda=100$ agree with the
exact primitive within $10^{-62}$. A linear tadpole subtraction has zero
second derivative: it does not remove this finite bubble contribution.

\Lc With canonical kinetic coefficient one, the low-frequency quadratic
effective action therefore has elastic modulus $1+\hbar_BW_g''(0)$ and
\begin{equation}\label{eq:sg}
 s_{1,g}=\tfrac12W_g''(0)=-\frac9{4\pi}.
\end{equation}
Each external potential fluctuation at a cubic vertex carries a spatial
derivative. The static bubble fixes the elastic $q^2$ term, and the uniform
kinetic mass density remains one. Equation~\eqref{eq:sg} is the
\emph{thermodynamic acoustic coefficient} inferred from that elasticity;
it is not an assertion that a perturbative on-shell phonon self-energy
has a uniform expansion at an acoustic threshold. This directly refutes
the audit's assertion $s_{1,\mathrm{cont}}=0$ for the stated bare-normalized
continuum calculation. Defining the \emph{renormalized} speed to be one
instead requires an additional finite subtraction, absent from that audit.

\subsection{Continued bad vacuum and lattice vacuum}
\Lc/\Ld (\ccheck) Continue the \emph{same prescription} as
$F_b(k)=-iF_g(ik)$, rather than treating unstable modes as a real vacuum.
The local density rule is
\begin{equation}\label{eq:Wb}
 W_{b,\rm ac}(U)=-W_g(-U)
  =\frac{(1+6U)^{3/2}-1-9U}{6\pi},\qquad
 s_{1,b,\rm ac}=\frac9{4\pi}.
\end{equation}
The sign is fixed by the spatial continuation as well as $U_g=-U_b$:
for the branch $k_g=ik_b$ one has $x_g=-ix_b$, and
$-i\int\dd x_g W_g(U_g)=-\int\dd x_b W_g(-U_b)$.
Equivalently, retain the real analytic finite part of the integrated
$q\sqrt{1+6U-q^2}$ primitive, then subtract its constant and linear terms.
The algebra is (c); interpreting this as a physical bad-equation vacuum
is (d). A fixed real stable-band cutoff is a different prescription and
does not justify this continuation.

\Lc (\ccheck) The Toda counterpart is real and requires no continuation.
Keep the reference $b$ and coordinates of \eqref{eq:dictionary} fixed.
A homogeneous Boussinesq strain changes lattice spacing by
$r-a=-6U$, so the harmonic frequency in \eqref{eq:todavac} is multiplied
by $e^{3U}$. The conversion of energy density is $1/(36b^2)$ and that
of $h$ is $3b$. Subtracting constant and linear terms gives
\begin{equation}\label{eq:WT}
 W_T^{B}(U)=\frac{e^{3U}-1-3U}{6\pi},\qquad
 \tfrac12(W_T^{B})''(0)=\frac3{4\pi}.
\end{equation}
The linear subtraction fixes background pressure/topological energy; it
does not affect the sound coefficient. This is the continuum local limit
of the \emph{lattice's own harmonic mode sum}, not a fitted soliton number.

\subsection{Matching the elastic modulus}
\Lc (\ccheck) Equations~\eqref{eq:Wb} and \eqref{eq:WT} determine
\begin{equation}\label{eq:match}
 W_T^B-W_{b,\rm ac}
 =\frac{e^{3U}-(1+6U)^{3/2}+6U}{6\pi}
 =-\frac3{2\pi}U^2+O(U^3).
\end{equation}
Here is the precise prescription required only to this order in amplitude:
\begin{enumerate}
\item Keep the conditional equal-mode-count, exactly tadpole-subtracted
trace and its continuation \eqref{eq:F}, with the stated momentum and energy
origins.
\item Match the vacuum elastic modulus in the same reference coordinates
to the microscopic Toda harmonic mode sum. Equivalently add
\begin{equation}\label{eq:counterterm}
 \boxed{\Delta\mathcal H=\hbar_B\lambda u^2,
       \qquad\lambda=s_{1,T}-s_{1,b,\rm ac}=-\frac3{2\pi}.}
\end{equation}
No order-$\hbar_B$ kinetic or canonical-momentum rescaling is introduced.
\item Evaluate this term on the classical soliton when computing the
first-order energy at fixed canonical $P$. Profile changes do not contribute
at this order, since the unperturbed profile is stationary for $H-vP$.
\end{enumerate}
\Lc (\ccheck) By \eqref{eq:moments} the resulting trace is
\begin{equation}\label{eq:matched}
 F_{\rm matched}(k)=F_{\rm ac}(k)-\frac{k^3}{\pi},\qquad
 \boxed{\lim_{k\to0}\frac{F_{\rm matched}(k)}{P_0(k)}=\frac1{2\pi}.}
\end{equation}
The equality is exact for the \emph{limiting coefficient}. We do not assert
that this function at finite $k$ equals Toda's dispersion. The full local
difference in \eqref{eq:match}, as well as derivative matching terms, would
require higher-amplitude matching; $\int u_0^3=O(k^5)$ already shows why
it cannot affect the present $k^3$ coefficient.

\paragraph{Natural or tuned?}
\Lc The coefficient in \eqref{eq:counterterm} is forced if ``lattice
compatible'' means matching the quantum vacuum elasticity of the specified
Toda Hamiltonian, at fixed mean strain and in the inherited units. It was
computed from the two vacuum mode sums, without using either soliton
coefficient. This is natural finite matching of a low-energy parameter.
It is not forced by the truncated Boussinesq equation alone. If one simply
postulated this counterterm without its vacuum condition, it would be a
tuned finite subtraction.

\Lc The distinction is material. The lattice classical local potential
in these units is
$(e^{6U}-1-6U)/36=U^2/2+U^3+3U^4/2+\cdots$.
Terms absent from the cubic continuum truncation and the full lattice
dispersion participate in its quantum elasticity. Calling
\eqref{eq:counterterm} a different \emph{regularisation scheme} means
matching these finite UV contributions into the effective Hamiltonian;
it is not a claim that two fixed bare Hamiltonians have identical physical
sound velocities. With renormalized parameters held fixed, a regulator
change alone cannot change a physical prediction.

\Lc The diagnostic cancellation is particularly transparent:
\begin{equation}\label{eq:soundsubtract}
 C_{\rm ac}-s_{1,b,\rm ac}
 =\frac2\pi-\frac9{4\pi}=-\frac1{4\pi}
 =C_T-s_{1,T}.
\end{equation}
The audit subtracted Toda's sound shift but set the continuum one to zero.
Its quoted ratio $-8$ therefore compares unequal acoustic baselines.

\section{Other candidate prescriptions and the good equation}
\paragraph{Cutoff endpoint and finite tadpole.}
\Lc (\ccheck) In the good trace the required equal-mode-count endpoint
has $B_\Lambda=3k\Lambda/(2\pi)+O(\Lambda^{-1})$ and no finite constant
term. Omitting it with the same tadpole leaves a divergence, as the earlier
determinant audit found. A finite redefinition of that endpoint can encode
an arbitrary finite subtraction, but no alternative endpoint quantization
was derived here from the lattice. In particular merely shifting a cutoff
is not an independent derivation of \eqref{eq:counterterm}.

\Lc A finite change to the Wick contraction of the cubic interaction
adds a constant multiple of $M=2k$. Since $P\sim2k^3/3$, a nonzero such
term creates a $k^{-2}$ divergence in $F/P$. It cannot convert one finite
coefficient into the other. Once the finite soft limit is imposed, that
linear freedom is fixed. The $u^2$ matching term is different: it changes
the elastic modulus and survives the old origin/mass annihilator
$\partial_\beta^3-\partial_\beta$ with limiting coefficient $\lambda$.

\paragraph{Real stable band (Scheme A).}
\Lc (\ccheck) Fresh quadrature of the exact Hirota transmission, using
$q=\sin t$ and an independently checked rational integrand in $z=\cos t$,
confirms the stipulated continuous-half-band trace
\begin{equation}\label{eq:band}
 F_{\rm band}=F_{\rm ac}-\frac{\sqrt3}{2}(k^2+\tfrac34),\qquad
 F_{\rm band}/P_0\sim-\frac{9\sqrt3}{16k^3}.
\end{equation}
The correct bad Wick term is $-3k/\pi$. Three 70-digit checks give errors
below $10^{-60}$. Subtracting the threshold constant alone still leaves an
$O(k^2)$ obstruction. The elastic counterterm \eqref{eq:counterterm}
does not remove either term. Consequently this specified Scheme A does
\emph{not} match Toda; projecting out further modes would require a new
spectral domain and boundary prescription, not an unexplained deletion
of \eqref{eq:band}'s extra term.

\paragraph{What the lattice mode sum has and has not supplied.}
\Lc It has supplied \eqref{eq:WT}, and hence a non-tuned leading local
matching condition. The independent Toda edge calculation supplies the
remaining $-1/(4\pi)$, which agrees with the continued continuum trace
after its own sound subtraction. \Ld We have not constructed the lattice
soliton fluctuation determinant and taken its complete continuum limit,
or proved that its constrained measure and endpoints produce
\eqref{eq:matched} at finite amplitude. Thus the present result is a
solution for the requested leading coefficient, not a full regulator
equivalence theorem.

\paragraph{Good equation.}
\Lc (\ccheck) The well-posed good vacuum needs no unstable-band contour
choice. Its finite elastic bubble is UV convergent after the stated linear
subtraction; the prescribed value $s_{1,g}=-9/(4\pi)$ is not a UV divergence
to discard. For its conditional soliton trace,
\begin{equation}
 C_g=-\frac2\pi,\qquad C_g-s_{1,g}=+\frac1{4\pi}.
\end{equation}
Absence of an unstable band does not preclude a finite $\hbar_B\lambda_g u^2$
renormalisation: it shifts both $s_{1,g}$ and $C_g$ by $\lambda_g$.
If the physical sound speed is defined to remain one, the corresponding
finite coefficient is $\lambda_g=9/(4\pi)$ and the soliton slope relative
to that sound line is $+1/(4\pi)$. With the bare Hamiltonian and the
normal-ordering condition fixed, no such extra subtraction is required.
\Ld Toda in the present real scaling regulates the bad branch, so it
does not force the same correction on the good equation. Neither a common
finite-amplitude quantum completion nor the remaining good projected
determinant domain follows from this calculation. ``Scheme independent''
is justified only after fixing the acoustic renormalisation and boundary
conditions, not from linear well-posedness alone.

\section{Verdict, reproducibility and failures recorded}
\begin{longtable}{@{}p{.37\linewidth}p{.43\linewidth}p{.13\linewidth}@{}}
\toprule Audited claim & Assessment & Check\\\midrule\endhead
Failed $3\,10^{-6}$ assertion & \Lc Incorrect numerical accuracy claim;
limiting identity confirmed. & edge\\
Six dictionary factors & \Lc Confirmed at leading classical order;
do not exclude finite quantum matching. & dictionary\\
$C_T=3/(4\pi)-1/(4\pi)$ & \Lc Confirmed for the nested edge limit;
uniform fixed-$P$ interchange unproved (d). & all three\\
$C_{\rm ac}=2/\pi$ & \Lc Confirmed as conditional continued trace;
physical determinant identification (d). & continuum\\
Light-root cancellation / all short-wave origin & \Lc Log weight vanishes;
literal spectral-localization inference is unsupported (d). & continuum\\
Continuum sound unrenormalised & \Lc Refuted: good $-9/(4\pi)$,
continued bad $+9/(4\pi)$. & continuum\\
Scheme A equals Toda after a tadpole change & \Lc Refuted for the stated
trace; nonzero even threshold term remains. & continuum\\
Matching prescription & \Lc $\Delta\mathcal H=-3\hbar_Bu^2/(2\pi)$,
fixed by lattice vacuum elasticity, gives $C=1/(2\pi)$. & continuum\\
Genuinely physical factor four established & \Lc Rejected as an inference:
the acoustic baselines were not matched. Full quantum equivalence (d). & all three\\
\bottomrule
\end{longtable}
The check column abbreviates script names: add the prefix
\texttt{onebody\_scheme\_} and the suffix \texttt{.py}.
Each has a companion JSON with values, precision, and relevant limitations.
Every formula newly claimed above
as (c) is checked by symbolic algebra, independent quadrature, or the
numerical asymptotics in these scripts; the stated physical and uniform-limit
qualifications are not promoted to numerically proved theorems.

\paragraph{Conclusion.}
\Lc Choose \textbf{(A) for the factor-of-four coefficient}: the
lattice-compatible prescription \eqref{eq:counterterm} is identified and
reproduces $C_T$ exactly, conditional on the inherited continued trace.
It is a finite vacuum-elastic matching prescription, not an arbitrary
rescaling of $\hbar$ or a mass tadpole. \Ld If (A) is instead meant to
assert equality of fully defined continuum and lattice quantum theories,
that stronger question is undetermined: the physical bad contour/domain,
boundary-compatible projected determinant, and a uniform quantum matching
construction are still missing. No claim of (B) follows from the audited
calculations.

\paragraph{Failures recorded.}
\begin{enumerate}
\item The unchanged original audit failed at line 383. The measured maximum
residual is $2.2546734882\,10^{-4}$; three rows ($y=16,32,64$) fail its
$3\,10^{-6}$ bound. Its statement of agreement below $2.8\,10^{-6}$ is
false for the files inspected. Nothing in that file was changed.
\item The fresh composite sea with 16-point panels at $h=.02$ failed the
endpoint-product tolerance $2\,10^{-11}$ with residual
$3.3117711365\,10^{-8}$. A 24-point retry also failed, at
$4.4067346994\,10^{-11}$. The final calculation uses 32-point panels
and a 40-point refinement; the original tolerance is retained.
\item The finer cubic fit still fails $3\,10^{-6}$ at $y=64$.
Quartic and quintic fits are explicitly reported separately; a successful
higher-order fit is not relabelled as success of the cubic or old fit.
\item The audit's $b=\pi/a$, intermediate velocity subtraction, energy
$e_0/y^2$ tail, and resonance kinematics fail the independent algebra as
detailed above. The zero continuum sound shift fails the vacuum curvature
calculation. The broad short-wave and physical-discrepancy conclusions
are therefore qualified rather than inherited.
\item The specified stable-band trace has a divergent soft ratio and is
not repaired by the local elastic matching term. We retain that failed
matching candidate instead of removing its even threshold contribution.
\item The initial PDF build reported two overfull lines. The transform
display and reproduction paragraph were reflowed; the final two-pass build
has no overfull lines or unresolved references. An auxiliary log-check
assertion initially mistook the package description ``Rerun checks for
auxiliary files'' for a rerun warning. Inspection identified that false
positive; the corrected check searches warning records instead.
\end{enumerate}

\paragraph{Reproduction and scope of execution.}
Run the three new Python scripts from their directory, then run
\texttt{pdflatex -interaction=nonstopmode -halt-on-error onebody\_scheme.tex}
twice in \texttt{phase2/}. Every new script ends with
\texttt{ALL CHECKS PASSED}. The original audit's expected failure is recorded
separately. No repository-wide historical sweep was run: existing scripts
can write non-prefix outputs, forbidden by this task. Existing input data
were read only; the edge script records and rechecks their SHA-256 digests.

\section{Internal sources read}\label{sec:sources}
No external literature is used as an independently verified authority here.
The internal documents used are \texttt{audit\_onebody.tex},
\texttt{audit\_det\_projected.tex}, \texttt{det\_projected\_oneloop.tex},
\texttt{onebody\_dispersion\_test.tex}, \texttt{onebody\_toda\_precision.tex},
\texttt{toda\_route.tex}, and relevant sections of \texttt{oneloop.tex}
and \texttt{frame\_dictionary.tex}. The old audit, one-body and frame
analysis scripts, Toda classical/scaling routines, and precision edge/sea
implementations were read to locate conventions and tests. Their
mathematical inputs are rechecked above without importing their code.
Saved finite-sea data are identified explicitly where reused; fresh
calculations are distinguished from saved-data reanalysis.
\end{document}
